\section{Mathematical model}

The model with time coordinate $t$, space coordinate $x$ and
main variables $T$ (temperature), $S_{\text{o}}$, $Y$ (oxygen
fraction) is composed by the equations describing the heat
balance~\eqref{Ec1a}, molar balance for oxygen,~\eqref{Ec2a} and
molar balance for oil~\eqref{Ec3a}:
\begin{equation}
	C_{m}\diffp{T}{t}+
	\diffp*{[c_{\text{o}}\rho_{\text{o}}S_{\text{o}}u_{\text{o}}\left(T-T_{\text{res}}\right)]}{x}=
	\rho_{\text{o}}Q_{\text{r}}YS_\text{g}S_{\text{o}}W_{\text{r}}
	\textcolor{red}{-\alpha\left(T-T_{\text{res}}\right)}.
	\label{Ec1a}
\end{equation}
\begin{equation}
	\varphi
	\diffp*{\left(YS_\text{g}\rho_{\text{g}}\right)}{t}+
	\diffp*{\left(YS_\text{g}\rho_{g}u_{\text{g}}\right)}{x}=
	-w\varphi
	\rho_{\text{o}}YS_\text{g}S_{\text{o}}W_{\text{r}}.
	\label{Ec2a}
\end{equation}
\begin{equation}
	\varphi
	\diffp*{\left(\rho_{\text{o}}S_{\text{o}}\right)}{t}+
	\diffp*{\left(\rho_{\text{o}}S_{\text{o}}u_{\text{o}}\right)}{x}=
	-\varphi\rho_{\text{o}}YS_\text{g}S_{\text{o}}W_{\text{r}}.
	\label{Ec3a}
\end{equation}
Where $C_{m}$ is the heat capacity of the matrix, $c_{\text{o}}$ is
the oil molar heat capacity, $\frac{\rho_{\text{o}}}{\rho_{\text{g}}}$
are oil/gas densities, $u_{\text{o}}$/$u_{\text{g}}$ are oil/gas
velocities, $T_{\text{res}}$ is the reservoir initial temperature,
$Q_{\text{r}}$ is the oil reaction enthalpy, $\varphi$ is the
porosity, and $\omega$ is a stochiometric coefficient for the
oil-oxygen reaction.
Following \cite{traveling,gargar2020combustion}, the reaction rate
$W_{\text{r}}$ is given by Arrhenius law:
\begin{equation}
	W_{\text{r}}=
	k_{p}\exp
	\Big(-\frac{E_\text{r}}{RT}\Big),
	\label{Ec4a}
\end{equation}
where $E_\text{r}$ is the activation energy, $R$ is the gas constant,
and $k_{p}$ is the pre-exponential coefficient.

In order to allow rigorous mathematical analysis, we follow
\cite{traveling,chapiro2012asymptotic,chapiro2014combustion,zavala2020classification,chapiro2015combustion} and assume constant oil and gas partial velocities $u_o$ and $u_g$.
As our goal is to compare our results to a more realistic model and experimental data.
%, we explain the best way to estimate these velocities in Appendix \ref{sec:appendix_velocities}.

\subsection{Dimensionless model}

In order to allow the mathematical analysis of the system \eqref{Ec1a}-\eqref{Ec3a}, we change the variables to dimensionless ones:
\begin{equation}\label{eq5}
	\bar{t} = \frac{t}{t^{\ast}},
	\;\;
	\bar{x} = \frac{x}{x^{\ast}},
	\;\;
	\bar{\theta} = \frac{T - T_{\text{res}}}{T^{\ast}},
	\;\;
	T^{\ast} = T_{\text{res}},
	\;\;
	t^{\ast} = \frac{10^{8}}{k_{p}},
	\;\;
	x^{\ast} = \frac{\rho_g}{\rho^{res}}\frac{u}{k_{p}},
\end{equation}
where bars indicate the dimensionless variables and $t^{\ast}$, $x^{\ast}$ are reference values for time and space.

Using \eqref{eq5} and omitting the bars, the system given by \eqref{Ec1a}--\eqref{Ec3a} is written in dimensionless form as:
\begin{equation}
	\label{eq:14}
	\displaystyle\frac{\partial \theta}{\partial t} + a_1 \displaystyle\frac{\partial }{\partial x}(S_{\text{o}} \theta) = b_1 S_{\text{o}} S_{\text{y}} \Phi - \textcolor{red}{\beta \theta},
\end{equation}
\begin{equation}
	\label{eq:15}
	\displaystyle\frac{\partial S_{\text{y}}}{\partial t} + a_2 \displaystyle\frac{\partial S_{\text{y}}}{\partial x} = - b_2 S_{\text{o}} S_{\text{y}} \Phi,
\end{equation}
\begin{equation}
	\label{eq:16}
	\displaystyle\frac{\partial S_{\text{o}}}{\partial t} + a_3 \displaystyle\frac{\partial S_{\text{o}}}{\partial  x} =
	- b_3 S_{\text{o}} S_{\text{y}} \Phi,
\end{equation}
where we introduced the variable $S_{\text{y}}=S_\text{g} Y$, $\Phi = t^{*} W_{\text{r}}$, and the constants are given by:
\begin{equation}\label{eq:17}
	a_1 = \displaystyle\frac{t^{*} u_o c_o \rho_o }{C_m x^{*}}, \;
	a_2 = \displaystyle\frac{t^{\ast} u_g}{\varphi x^{\ast}}, \;
	a_3 = \displaystyle\frac{t^{\ast} u_o}{\varphi x^{\ast}}, \;
	b_1 = \displaystyle\frac{\rho_o  Q_r }{C_m T^{*}}, \;
	b_2 = \displaystyle\frac{w \rho_o}{\rho_g},\;
	b_3 = 1, \\
	\textcolor{red}{\beta = \displaystyle\frac{\alpha t^{*}}{C_m}}
\end{equation}

In this paper we study the combustion front and
thus we assume that the reaction does not occur at the
boundaries, i.e., the reaction rate in \eqref{eq:14}-\eqref{eq:16} vanishes when: there
is no oxygen ($S_{\text{y}} = 0$), which we call ``Oxygen Controlled'' case (OC);
there is no fuel ($S_{\text{o}} = 0$), which we call ``Fuel Controlled'' case (FC).

In order to approximate applications, we consider that the left in-
jection end of the domain is FC, and the production right end
of the domain is OC. Because of the non-dimensionalization pre-
sented above we have that $S_{\text{o}} = 1$ in OC state and $S_{\text{y}} = 1$ in FC state.

Thus, the investigated domain corresponds to non-negative satura-
tions $S_{\text{o}} \geq 0$ and $S_{\text{y}} \geq 0$ and is restricted to $0<x<\infty$, $t\geq 0$, with
the (constant) boundary conditions:

\subsection{Initial and boundary conditions}
To complete system \eqref{eq:14}-\eqref{eq:16}, we need initial and boundary conditions. Physically, there is no chemical reaction in the reservoir before the combustion and no reaction happening in the injection well. Mathematically, this means that the source terms in \eqref{eq:14}-\eqref{eq:16} are equal to zero ($S_{\text{y}} S_{\text{o}} \Phi = 0$), as the temperature is always positive, this equation yields $S_{\text{y}} = 0$ or $S_{\text{o}} = 0$. Thus, we can expect the following options for the limit states (initial and injection conditions).
\begin{enumerate}
	\item If $S_{\text{y}}= S_\text{g} Y = 0$ then  $Y = 0$ or $S_\text{g} = 0$, therefore:
	      \begin{enumerate}
		      \item If $Y = 0$ then $S_{\text{y}} = 0$ and the limit state is:
		            \begin{equation}\label{eq8c}
			            (\theta,S_{\text{y}},S_{\text{o}}) = (\theta,0,S_{\text{o}}).
		            \end{equation}
		      \item If $S_\text{g} = 0$ then $S_{\text{o}} = 1$, $S_{\text{y}} = 0$, and the limit state is:
		            \begin{equation}\label{eq9c}
			            (\theta,S_{\text{y}},S_{\text{o}}) = (\theta,0,1).
		            \end{equation}
	      \end{enumerate}
	\item If $S_{\text{o}} = 0$ then $S_\text{g}=1$, and the limit state is:
	      \begin{equation}\label{eq10c}
		      (\theta,S_{\text{y}},S_{\text{o}}) = (\theta,S_{\text{y}},0).
	      \end{equation}
\end{enumerate}
In summary, we have two cases for boundary conditions: \eqref{eq8c} and \eqref{eq9c} with no oxygen, and \eqref{eq10c} without oil.
Based on the physical problem we are interested in, we consider no oil in the injected fluid (left condition) and no oxygen in the initial (right) conditions.

In the present paper, we study the Riemann problem composed of
equations \eqref{eq:14}-\eqref{eq:16} with initial conditions defined
as a step function with no oil as the left condition ($x<0$) and no
oxygen as the right condition ($x>0$):
\begin{equation}
	\label{eq:Riemann_boundary}
	(\theta^{\text{L}},S_{\text{y}}^{\text{L}},S_{\text{o}}^{\text{L}})=
	\left(0,S_{\text{y}}^L,0\right)
	\text{ and }
	\left(\theta^{\text{R}},S_{\text{y}}^{\text{R}},S_{\text{o}}^{\text{R}}\right)=
	\left(0,0,S_{\text{o}}^{\text{R}}\right),
\end{equation}
where it was assumed that the injected air was not heated.
Thus, both initial and injection temperatures coincide with the reservoir temperature and are equal to zero following \eqref{eq5}.
\cite{zavala2020classification}
